%%%%%%%%%%%%%%%%%%%%%%%%%%%%%%%%%%%%%%%%%%%%%%%%%%%%%%%%%%%%%%%%%%%%%%
% Lớp 8 | Chương IV. Định lí Thalès | Định lí Thalès trong tam giác
% Dạng 8H4?15-4: Ứng dụng thực tế (đo chiều cao, khoảng cách)
%%%%%%%%%%%%%%%%%%%%%%%%%%%%%%%%%%%%%%%%%%%%%%%%%%%%%%%%%%%%%%%%%%%%%%


% nguồn: dan_tu_AI (nhập 16/11/2026)
\begin{bt}%[8H4H15-4]%[ma:82fc8e62]
	Một cái cọc $DE$ cao $1{,}5\text{ m}$ được cắm thẳng đứng trên mặt đất. Người ta thấy đầu bóng của cọc trùng với đầu bóng của một cái cây $BC$ (cây và cọc cùng vuông góc với mặt đất, chân cọc $D$ nằm giữa đầu bóng $A$ và gốc cây $B$). Đo được $AD=2\text{ m}$, $AB=10\text{ m}$. Tính chiều cao $BC$ của cây.
	\loigiai{
		Cọc $DE$ và cây $BC$ cùng vuông góc với mặt đất nên $DE\parallel BC$.\\
		Tia nắng đi qua $A$, $E$, $C$ nên $A$, $E$, $C$ thẳng hàng. Trong tam giác $ABC$ có $DE\parallel BC$, theo hệ quả của định lí Thalès:
		\[ \dfrac{DE}{BC}=\dfrac{AD}{AB}\Rightarrow BC=\dfrac{DE\cdot AB}{AD}=\dfrac{1{,}5\cdot 10}{2}=7{,}5\text{ (m)}. \]
		Vậy cây cao $7{,}5\text{ m}$.
	}
\end{bt}

% nguồn: dan_tu_AI (nhập 16/11/2026)
\begin{bt}%[8H4H15-4]%[ma:0ab6d248]
	Để đo khoảng cách từ điểm $B$ trên bờ sông đến điểm $A$ ở bờ bên kia, người ta làm như sau: chọn điểm $D$ sao cho $A$, $B$, $D$ thẳng hàng và $BD=10\text{ m}$; chọn điểm $C$ và điểm $E$ sao cho $A$, $C$, $E$ thẳng hàng và $BC\parallel DE$. Đo được $BC=30\text{ m}$, $DE=40\text{ m}$. Tính khoảng cách $AB$.
	\loigiai{
		Trong tam giác $ADE$ có $BC\parallel DE$ ($B\in AD$, $C\in AE$). Theo hệ quả của định lí Thalès:
		\[ \dfrac{AB}{AD}=\dfrac{BC}{DE}=\dfrac{30}{40}=\dfrac34. \]
		Mà $AD=AB+10$ nên $\dfrac{AB}{AB+10}=\dfrac34$, suy ra $4AB=3AB+30$, do đó $AB=30\text{ (m)}$.\\
		Vậy khoảng cách $AB$ là $30\text{ m}$.
	}
\end{bt}

% nguồn: dan_tu_AI (nhập 16/11/2026)
\begin{bt}%[8H4H15-4]%[ma:f765f902]
	Một cột điện $BC$ thẳng đứng được giữ bởi sợi dây néo căng thẳng $AC$ ($A$ trên mặt đất). Người ta cắm một cọc $DE$ thẳng đứng sao cho chân cọc $D$ nằm trên đoạn $AB$ và đầu cọc $E$ chạm vào dây néo. Biết $AD=1{,}2\text{ m}$, $DE=0{,}8\text{ m}$, $AB=6\text{ m}$. Tính chiều cao của cột điện.
	\loigiai{
		Cọc $DE$ và cột $BC$ cùng vuông góc với mặt đất nên $DE\parallel BC$.\\
		Trong tam giác $ABC$ có $DE\parallel BC$, theo hệ quả của định lí Thalès:
		\[ \dfrac{AD}{AB}=\dfrac{DE}{BC}\Rightarrow BC=\dfrac{DE\cdot AB}{AD}=\dfrac{0{,}8\cdot 6}{1{,}2}=4\text{ (m)}. \]
		Vậy cột điện cao $4\text{ m}$.
	}
\end{bt}

% nguồn: dan_tu_AI (nhập 16/11/2026)
\begin{bt}%[8H4H15-4]%[ma:e66b1e99]
	Hai điểm $A$ và $B$ ở hai phía của một cái hồ nên không đo trực tiếp được khoảng cách. Người ta chọn một điểm $O$ ở ngoài hồ, lấy điểm $C$ trên đoạn $OA$ và điểm $D$ trên đoạn $OB$ sao cho $OC=\dfrac13OA$, $OD=\dfrac13OB$. Đo được $CD=25\text{ m}$. Tính khoảng cách $AB$.
	\loigiai{
		Ta có $\dfrac{OC}{OA}=\dfrac{OD}{OB}=\dfrac13$ nên theo định lí Thalès đảo, $CD\parallel AB$.\\
		Theo hệ quả của định lí Thalès trong tam giác $OAB$:
		\[ \dfrac{CD}{AB}=\dfrac{OC}{OA}=\dfrac13\Rightarrow AB=3\cdot CD=3\cdot 25=75\text{ (m)}. \]
	}
\end{bt}

% nguồn: dan_tu_AI (nhập 16/11/2026)
\begin{bt}%[8H4H15-4]%[ma:cba5e45a]
	Bạn An có mắt cách mặt đất $1{,}5\text{ m}$ muốn đo chiều cao của một cái cây. An cắm một cọc thẳng đứng cao $2{,}5\text{ m}$ cách chỗ mình đứng $2\text{ m}$, rồi lùi lại sao cho mắt, đỉnh cọc và ngọn cây thẳng hàng; khi đó An đứng cách gốc cây $8\text{ m}$ (chân An, chân cọc và gốc cây thẳng hàng). Tính chiều cao của cây.
	\loigiai{
		Gọi $E$ là mắt An. Kẻ đường thẳng qua $E$ song song với mặt đất, cắt cọc tại $K$ và cắt cây tại $H$. Gọi $P$ là đỉnh cọc, $Q$ là ngọn cây.\\
		Khi đó $EK=2\text{ m}$, $EH=8\text{ m}$, $KP=2{,}5-1{,}5=1\text{ (m)}$.\\
		Cọc và cây cùng thẳng đứng nên $KP\parallel HQ$. Trong tam giác $EHQ$, theo hệ quả của định lí Thalès:
		\[ \dfrac{KP}{HQ}=\dfrac{EK}{EH}=\dfrac28\Rightarrow HQ=4\cdot KP=4\text{ (m)}. \]
		Chiều cao của cây: $4+1{,}5=5{,}5\text{ (m)}$.
	}
\end{bt}

% nguồn: dan_tu_AI (nhập 16/11/2026)
\begin{bt}%[8H4V15-4]%[ma:0be4ff6b]
	Để đo chiều cao $AH$ của một tòa nhà, người ta cắm hai cọc thẳng đứng $CD=2\text{ m}$ và $EF=3\text{ m}$ ($D$, $F$ là chân cọc) sao cho các điểm $O$, $D$, $F$, $H$ thẳng hàng trên mặt đất và các điểm $O$, $C$, $E$, $A$ thẳng hàng. Biết khoảng cách giữa hai cọc $DF=4\text{ m}$ và khoảng cách từ cọc $EF$ đến chân tòa nhà là $FH=36\text{ m}$. Tính chiều cao của tòa nhà.
	\loigiai{
		Hai cọc và tòa nhà cùng vuông góc với mặt đất nên $CD\parallel EF\parallel AH$.\\
		Tam giác $OFE$ có $CD\parallel EF$, theo hệ quả của định lí Thalès:
		\[ \dfrac{OD}{OF}=\dfrac{CD}{EF}=\dfrac23\Rightarrow\dfrac{OD}{OD+4}=\dfrac23\Rightarrow 3OD=2OD+8\Rightarrow OD=8\text{ (m)}. \]
		Suy ra $OF=12\text{ m}$, $OH=OF+FH=12+36=48\text{ (m)}$.\\
		Tam giác $OHA$ có $EF\parallel AH$, theo hệ quả của định lí Thalès:
		\[ \dfrac{EF}{AH}=\dfrac{OF}{OH}=\dfrac{12}{48}=\dfrac14\Rightarrow AH=4\cdot 3=12\text{ (m)}. \]
		Vậy tòa nhà cao $12\text{ m}$.
	}
\end{bt}

% nguồn: dan_tu_AI (nhập 18/11/2026)
\begin{ex}%[8H4V15-4]%[ma:b48e57ab]
	Bóng của một cột cờ trên mặt đất dài $6\text{ m}$. Cùng thời điểm đó, một thanh sắt cao $1{,}5\text{ m}$ cắm vuông góc với mặt đất có bóng dài $1\text{ m}$. Chiều cao cột cờ là
	\choice
	{$4\text{ m}$}
	{$6\text{ m}$}
	{$8\text{ m}$}
	{\True $9\text{ m}$}
	\loigiai{
		Gọi chiều cao cột cờ là $h$. Vì các tia nắng song song, ta có tỉ lệ: $\dfrac{h}{1{,}5} = \dfrac{6}{1} \Rightarrow h = 9\text{ m}$.
	}
\end{ex}

% nguồn: dan_tu_AI (nhập 18/11/2026)
\begin{ex}%[8H4C15-4]%[ma:9dc79322]
	Để đo bề rộng $AB$ của một khúc sông, một người đo được các khoảng cách $AC = 3\text{ m}$, $CD = 2\text{ m}$ (với $A, C, D$ thẳng hàng và vuông góc với bờ sông). Kẻ đường $DE \parallel \text{bờ sông}$ cắt tia nhìn từ $B$ qua $C$ tại $E$. Biết $DE = 4\text{ m}$, bề rộng khúc sông $AB$ là
	\choice
	{$5\text{ m}$}
	{\True $6\text{ m}$}
	{$8\text{ m}$}
	{$10\text{ m}$}
	\loigiai{
		Do $AB \parallel DE$ nên theo hệ quả Thalès: $\dfrac{AB}{DE} = \dfrac{AC}{CD} \Rightarrow \dfrac{AB}{4} = \dfrac{3}{2} \Rightarrow AB = 6\text{ m}$.
	}
\end{ex}

% nguồn: dan_tu_AI (nhập 18/11/2026)
\begin{ex}%[8H4N15-4]%[ma:efa09116]
	Một cái cây cao $5\text{ m}$ có bóng dài $2\text{ m}$. Cùng lúc đó, bóng của cột điện dài $4\text{ m}$.
	\choiceTF
	{\True Tỉ số bóng của cây và bóng cột điện là $\dfrac{1}{2}$}
	{Cột điện cao $8\text{ m}$}
	{\True Cột điện cao $10\text{ m}$}
	{Tỉ số chiều cao cây và cột điện là $\dfrac{2}{5}$}
	\loigiai{
		\begin{itemchoice}
			\itemch Đúng. $2/4 = 1/2$.
			\itemch Sai. Chiều cao $h = 5 \cdot \dfrac{4}{2} = 10\text{ m}$.
			\itemch Đúng. $h = 10\text{ m}$.
			\itemch Sai. Tỉ số chiều cao là $\dfrac{5}{10} = \dfrac{1}{2}$.
		\end{itemchoice}
	}
\end{ex}

% nguồn: dan_tu_AI (nhập 18/11/2026)
\begin{ex}%[8H4H15-4]%[ma:f47899e1]
	Một người cao $1{,}6\text{ m}$ đứng cách bờ sông $1\text{ m}$, nhìn thấy ảnh của ngọn cây (bờ bên kia) qua vũng nước cách chân người $2\text{ m}$. Cây cách vũng nước $10\text{ m}$. (Giả sử các tia sáng thẳng).
	\choiceTF
	{Tỉ số khoảng cách là $5$}
	{\True Cây cao $8\text{ m}$}
	{Khoảng cách từ người đến cây là $12\text{ m}$}
	{\True Nếu người lùi ra xa thêm, bóng cây vẫn giữ nguyên tỉ lệ với khoảng cách}
	\loigiai{
		\begin{itemchoice}
			\itemch Sai. Tỉ số khoảng cách đến vũng nước là $10/2 = 5$, nhưng ý nói chung chung không rõ. Coi như sai do không rõ ràng. (Thực ra $\triangle$ đồng dạng tỉ số 5).
			\itemch Đúng. $h/1{,}6 = 10/2 \Rightarrow h = 8$.
			\itemch Đúng. $2 + 10 = 12\text{ m}$.
			\itemch Đúng. Tỉ lệ ánh sáng phản xạ không đổi nếu chiều cao cây và người cố định.
		\end{itemchoice}
	}
\end{ex}

% nguồn: dan_tu_AI (nhập 18/11/2026)
\begin{ex}%[8H4C15-4]%[ma:0a1446f1]
	Hai cột đèn cao $a$ và $b$ cách nhau một khoảng $d$. Một người nhận thấy điểm giao nhau của hai bóng đèn trên mặt đất nằm ở độ cao $h$.
	\choiceTF
	{\True $\dfrac{h}{a} + \dfrac{h}{b} = 1$}
	{\True $h$ không phụ thuộc vào $d$}
	{Nếu $a=b$ thì $h=a$}
	{$h = \dfrac{a+b}{ab}$}
	\loigiai{
		\begin{itemchoice}
			\itemch Đúng. Áp dụng bài toán tương tự hình thang, $\dfrac{h}{a} + \dfrac{h}{b} = 1$.
			\itemch Đúng. Triệt tiêu khoảng cách $d$ trong tỉ lệ.
			\itemch Sai. Nếu $a=b$ thì $h = a/2$.
			\itemch Sai. $h = \dfrac{ab}{a+b}$.
		\end{itemchoice}
	}
\end{ex}

% nguồn: dan_tu_AI (nhập 18/11/2026)
\begin{bt}%[8H4H15-4]%[ma:657d47ea]
	Một tòa nhà có bóng dài $25\text{ m}$. Một cái cọc cao $2\text{ m}$ có bóng dài $1\text{ m}$. Tính chiều cao tòa nhà (m).
	\shortans{50}
	\loigiai{
		$h = 25 \cdot \dfrac{2}{1} = 50\text{ m}$.
	}
\end{bt}

% nguồn: dan_tu_AI (nhập 18/11/2026)
\begin{bt}%[8H4V15-4]%[ma:9e86d89e]
	Một chiếc thang $AB$ dài $5\text{ m}$ dựa vào tường. Khoảng cách từ chân thang $A$ đến tường là $3\text{ m}$. Một bậc thang $CD$ song song mặt đất cách chân thang $2\text{ m}$ dọc theo thang (tức $AC=2\text{ m}$). Tính khoảng cách từ $C$ đến tường (m).
	\shortans{1,8}
	\loigiai{
		Gọi điểm chạm tường là $B$, mặt đất là tường ở $H$. Khoảng cách cần tìm là $d$.
		Bằng Thalès: $\dfrac{BC}{BA} = \dfrac{d}{3}$.
		$BC = 5 - 2 = 3$. $\dfrac{3}{5} = \dfrac{d}{3} \Rightarrow d = 1{,}8\text{ m}$.
	}
\end{bt}

% nguồn: dan_tu_AI (nhập 18/11/2026)
\begin{bt}%[8H4C15-4]%[ma:e2a0fdfd]
	Để đo độ sâu của một cái giếng cạn hình trụ có đường kính $2\text{ m}$, người ta đặt mắt cách miệng giếng $0{,}5\text{ m}$ (theo phương ngang) và cách mặt đất $1{,}5\text{ m}$ thì vừa thấy được góc đáy giếng ở phía bên kia. Độ sâu của giếng là bao nhiêu m?
	\shortans{6}
	\loigiai{
		Dùng Thalès: $\dfrac{h_{\text{giếng}}}{1{,}5} = \dfrac{2}{0{,}5} = 4 \Rightarrow h = 4 \cdot 1{,}5 = 6\text{ m}$.
	}
\end{bt}

% nguồn: dan_tu_AI (nhập 18/11/2026)
\begin{bt}%[8H4N15-4]%[ma:b783d3d6]
	Một cột đèn cao $6\text{ m}$. Cùng thời điểm đó, bóng của một người cao $1{,}5\text{ m}$ trên mặt đất dài $2\text{ m}$. Tính chiều dài bóng của cột đèn.
	\loigiai{
		Vì các tia nắng mặt trời xem như song song, tạo ra các tam giác có cùng tỉ lệ các cạnh tương ứng.
		Gọi $x$ (m) là chiều dài bóng của cột đèn.
		Ta có tỉ lệ:
		\[ \dfrac{\text{Bóng cột}}{\text{Chiều cao cột}} = \dfrac{\text{Bóng người}}{\text{Chiều cao người}} \]
		\[ \Rightarrow \dfrac{x}{6} = \dfrac{2}{1{,}5} \Rightarrow x = 6 \cdot \dfrac{2}{1{,}5} = 8\text{ m}. \]
		Vậy bóng cột đèn dài $8\text{ m}$.
	}
\end{bt}

% nguồn: dan_tu_AI (nhập 18/11/2026)
\begin{bt}%[8H4H15-4]%[ma:1cc36353]
	Một người đo khoảng cách qua một con sông. Người đó đặt các cọc tại $A, B, C, D, E$ sao cho $A, B, C$ thẳng hàng; $A, D, E$ thẳng hàng và $BD \parallel CE$. Biết $AB = 10\text{ m}, BC = 15\text{ m}, BD = 8\text{ m}$. Tính bề rộng sông $CE$.
	\loigiai{
		Vì $BD \parallel CE$, áp dụng hệ quả định lí Thalès trong $\triangle ACE$:
		\[ \dfrac{AB}{AC} = \dfrac{BD}{CE} \]
		Ta có $AC = AB + BC = 10 + 15 = 25\text{ m}$.
		\[ \Rightarrow \dfrac{10}{25} = \dfrac{8}{CE} \Rightarrow CE = \dfrac{25 \cdot 8}{10} = 20\text{ m}. \]
	}
\end{bt}

% nguồn: dan_tu_AI (nhập 18/11/2026)
\begin{bt}%[8H4V15-4]%[ma:fdb88266]
	Một chiếc lều chữ A có các thanh chống như hình. Thanh ngang $MN$ song song với mặt đất $BC$. Biết độ dài thanh ngang là $1{,}5\text{ m}$, mặt đất $BC$ rộng $2\text{ m}$. Thanh chống dài $AM = 2{,}4\text{ m}$. Tính chiều dài từ $M$ đến chân lều $B$.
	\loigiai{
		Do $MN \parallel BC$, áp dụng hệ quả định lí Thalès trong $\triangle ABC$:
		\[ \dfrac{MN}{BC} = \dfrac{AM}{AB} \Rightarrow \dfrac{1{,}5}{2} = \dfrac{2{,}4}{AB} \]
		\[ \Rightarrow AB = \dfrac{2{,}4 \cdot 2}{1{,}5} = 3{,}2\text{ m}. \]
		Khoảng cách từ $M$ đến $B$ là:
		\[ MB = AB - AM = 3{,}2 - 2{,}4 = 0{,}8\text{ m}. \]
	}
\end{bt}

% nguồn: dan_tu_AI (nhập 18/11/2026)
\begin{bt}%[8H4C15-4]%[ma:5583f3ba]
	Một nhà thiết kế muốn tạo một cái thang xếp cắt chéo. Gọi trục giao nhau là $O$, đỉnh thang là $A, B$, chân là $C, D$ sao cho $AB \parallel CD$. Biết $AB = 30\text{ cm}, CD = 50\text{ cm}$. Nếu độ cao từ $O$ đến $AB$ là $h_1$, từ $O$ đến $CD$ là $h_2$. Tính tỉ số $\dfrac{h_1}{h_2}$.
	\loigiai{
		Kẻ đường cao chung đi qua $O$, vuông góc với $AB$ và $CD$.
		Theo định lí Thalès, mọi đoạn thẳng song song bị cắt bởi hai cát tuyến cắt nhau đều tỉ lệ.
		\[ \dfrac{h_1}{h_2} = \dfrac{AB}{CD} = \dfrac{30}{50} = \dfrac{3}{5}. \]
		Vậy tỉ số đường cao là $0{,}6$.
	}
\end{bt}

% nguồn: De_AI_Toan8_cuoi_hoc_ki_i_1 (nhập 21/11/2026)
\begin{bt}%[8H4H15-4]%[ma:65df5383]
	Để đo chiều cao $AB$ của một cột cờ (vuông góc với mặt đất tại $B$), người ta cắm một cái cọc $CD$ cao $1{,}5\text{ m}$ vuông góc với mặt đất tại $D$ sao cho bóng của đỉnh cột cờ $A$ và bóng của đỉnh cọc $C$ trùng nhau tại điểm $E$ trên mặt đất (biết ba điểm $B$, $D$, $E$ thẳng hàng và $D$ nằm giữa $B$, $E$). Biết khoảng cách $DE = 2\text{ m}$ và $BD = 10\text{ m}$. Tính chiều cao $AB$ của cột cờ.
	\loigiai{
		Vì cột cờ $AB$ và cọc $CD$ cùng vuông góc với mặt đất $BE$ nên $CD \parallel AB$. \\
Ta có $BE = BD + DE = 10 + 2 = 12\text{ (m)}$. \\
Xét $\triangle EAB$ có $CD \parallel AB$ ($C \in EA$, $D \in EB$). Theo hệ quả của định lí Thalès ta có: \\
$\dfrac{CD}{AB} = \dfrac{ED}{EB} \Rightarrow \dfrac{1{,}5}{AB} = \dfrac{2}{12} = \dfrac{1}{6}$. \\
Suy ra $AB = 1{,}5 \cdot 6 = 9\text{ (m)}$. \\
Vậy chiều cao của cột cờ là $9\text{ m}$.
	}
\end{bt}

% nguồn: De_AI_Toan8_cuoi_hoc_ki_i_1 (nhập 21/11/2026)
\begin{bt}%[8H4H15-4]%[ma:1330abb1]
	Để đo khoảng cách giữa hai địa điểm $A$ và $B$ ở hai bên bờ một hồ nước, người ta chọn một điểm $C$ trên đất liền sao cho có thể đo trực tiếp các đoạn thẳng $CA$ và $CB$. Trên đoạn thẳng $CA$ lấy điểm $M$ sao cho $CM = 20\text{ m}$, $MA = 60\text{ m}$; trên đoạn thẳng $CB$ lấy điểm $N$ sao cho $CN = 25\text{ m}$, $NB = 75\text{ m}$. Đo trực tiếp được khoảng cách $MN = 32\text{ m}$.
a) Chứng minh $MN \parallel AB$.
b) Tính khoảng cách giữa hai địa điểm $A$ và $B$.
	\loigiai{
		a) Ta có $CA = CM + MA = 20 + 60 = 80\text{ (m)}$ và $CB = CN + NB = 25 + 75 = 100\text{ (m)}$. \\
Xét $\triangle CAB$ có $\dfrac{CM}{CA} = \dfrac{20}{80} = \dfrac{1}{4}$ và $\dfrac{CN}{CB} = \dfrac{25}{100} = \dfrac{1}{4}$. \\
Suy ra $\dfrac{CM}{CA} = \dfrac{CN}{CB}$. Theo định lí Thalès đảo, ta có $MN \parallel AB$. \\
b) Vì $MN \parallel AB$ nên theo hệ quả của định lí Thalès trong $\triangle CAB$, ta có: \\
$\dfrac{MN}{AB} = \dfrac{CM}{CA} = \dfrac{1}{4} \Rightarrow AB = 4 \cdot MN = 4 \cdot 32 = 128\text{ (m)}$. \\
Vậy khoảng cách giữa hai địa điểm $A$ và $B$ là $128\text{ m}$.
	}
\end{bt}
